\textbf{Powerful Number Sieve}

\textbf{目标}
求积性函数前缀和
\[
F(n)=\sum_{i=1}^{n}f(i).
\]

\textbf{核心构造}
取积性函数 $g$，满足
\[
g(p)=f(p)
\]
且
\[
G(n)=\sum_{i=1}^{n}g(i)
\]
易求。
令 $h=f/g$ 表示 Dirichlet 卷积意义下的“商”，即
\[
f=g*h.
\]

\textbf{关键结论}
有
\[
h(1)=1,\qquad h(p)=0.
\]
因此 $h$ 仅可能在 Powerful Number 上非零。
若记 $\mathrm{PN}$ 为 Powerful Number 集合，则
\[
F(n)=\sum_{\substack{d\in \mathrm{PN}\\ d\le n}} h(d)\,G\!\left(\left\lfloor\frac{n}{d}\right\rfloor\right).
\]

\textbf{Powerful Number}
正整数 $n=\prod p_i^{e_i}$ 是 PN，当且仅当所有 $e_i\ge 2$；$1$ 也算 PN。

\textbf{PN 性质}
\begin{itemize}
\item 任意 PN 可写成 $a^2b^3$。
\item $[1,n]$ 内 PN 个数为 $O(\sqrt n)$。
\end{itemize}

\textbf{计算 $h(p^e)$}
由
\[
f(p^e)=\sum_{i=0}^{e}g(p^i)h(p^{e-i})
\]
得
\[
h(p^e)=f(p^e)-\sum_{i=1}^{e}g(p^i)h(p^{e-i}).
\]
若能直接写出 $h(p^e)$ 的闭式更好。

\textbf{快记}
PN 筛 = “把普通数的贡献折到 PN 上，再枚举 $O(\sqrt n)$ 个有效状态”。
